% glossaries.tex
% Acronym definitions for "Mind the Gap: A Structured Survey of AI Mental Health Technologies
% and Their Equity Failures for Marginalized Populations"
%
% Rules applied:
%   - AI, LLM, WHO, NEDA: abbreviation-only (no expanded definition body)
%   - ACM, APA, LGBTQ, LGBTQA: removed entirely
%   - No plural glossary entries; use \gls{llm}s / \gls{rct}s in prose

\newabbreviation{ai}{AI}{Artificial Intelligence}
\newabbreviation{llm}{LLM}{Large Language Model}
\newabbreviation{who}{WHO}{World Health Organization}
\newabbreviation{neda}{NEDA}{National Eating Disorders Association}
\newabbreviation{lgbtq}{LGBTQ+}{Lesbian, Gay, Bisexual, Transgender, and Queer}
\newabbreviation{rct}{RCT}{Randomized Controlled Trial}
\newabbreviation{adhd}{ADHD}{Attention Deficit Hyperactivity Disorder}
\newabbreviation{gad}{GAD-7}{Generalized Anxiety Disorder 7-item Scale}
\newabbreviation{fda}{FDA}{Food and Drug Administration}
\newabbreviation{cald}{CALD}{Culturally and Linguistically Diverse}

\newabbreviation[
  description={Cognitive Behavioural Therapy: A structured, evidence-based psychotherapeutic approach targeting the relationship between thoughts, feelings, and behaviours; the most common modality delivered by purpose-built mental health chatbots}
]{cbt}{CBT}{Cognitive Behavioural Therapy}

\newabbreviation[
  description={Digital Mental Health: The field encompassing mobile applications, online platforms, chatbots, and AI-powered tools designed to deliver mental health support outside traditional clinical settings}
]{dmh}{DMH}{Digital Mental Health}

\newabbreviation[
  description={Web Content Accessibility Guidelines: The internationally recognized technical standards for digital accessibility, published by the World Wide Web Consortium, specifying minimum requirements for screen reader compatibility, colour contrast, and interaction design}
]{wcag}{WCAG}{Web Content Accessibility Guidelines}
